\section{Limitations and Threats to Validity}
\label{sec:limitations}

\paragraph{Scope of the datasets.}
All datasets are small (at most 2{,}000 rows), fully numerical and complete. Our conclusions do not extend to larger
tables, categorical or text features, or missing values---settings that tabular foundation models are designed to
handle and in which their relative performance may differ. The datasets come from two curated suites; curation and the
selection filters shape which problems are represented. The filters relied on OpenML metadata, which was wrong for
\texttt{forest\_fires}: it was excluded at load time because of undeclared text columns, after results for 17 other
datasets existed. The exclusion follows the selection criterion and did not depend on any result, but it reduced the
regression suite from seven to six datasets, and we did not evaluate the dataset with its text columns encoded or
removed. Six regression datasets support descriptive conclusions only.

\paragraph{Dependence between datasets.}
Some datasets share a source (four NASA software-defect datasets; two \texttt{mfeat} feature sets of the same digits),
which violates the independence assumed by the Friedman and Wilcoxon tests. The planned family-collapsed analysis
addresses this only partly.

\paragraph{Untuned baselines and a single checkpoint.}
No model was tuned. Tuned gradient boosting or regularized linear models could narrow or reverse differences, and
library defaults are not designed for very small data. Only the \texttt{small} Kumo Tabular checkpoint was evaluated,
with float32 inference on an older GPU; the developers' claims concern larger checkpoints as well.

\paragraph{Pretraining overlap.}
The developers state that Kumo Tabular was pretrained on synthetic data only \citep{nvidia_kumo_blog}. We cannot verify
this; if any benchmark dataset or close variant was seen in pretraining, Kumo Tabular's performance on it would be
overstated.

\paragraph{Statistical power and multiplicity.}
With 15 classification datasets the cross-dataset tests can detect only consistent differences; absence of
significance is not evidence of equivalence. Secondary metrics and sensitivity analyses are not adjusted for
multiplicity and are descriptive. Within-dataset comparisons in the case studies rest on five dependent splits.

\paragraph{Timing of the analysis plan.}
The analysis plan was written after the complete results on \texttt{wdbc} had been seen, on which Kumo Tabular was best
on all quality metrics, and it is timestamped only in a local repository. Choices made in the plan, such as the
primary metrics, could in principle have been influenced by that one dataset. The analysis without \texttt{wdbc}
gives the same conclusions (Table~\ref{tab:sensitivity}), and the secondary metrics are reported in full.

\paragraph{Runtime.}
Runtimes come from one laptop with a GPU for Kumo Tabular and a CPU for all other models, and timings of a few
milliseconds vary between runs. They indicate orders of magnitude on this hardware, not algorithmic costs.

\paragraph{Construct validity.}
Accuracy and macro-F1 use the arg-max decision rule; in applications the decision threshold would be chosen for a
target cost. The medical case-study datasets are used as benchmarks only; no conclusion about clinical use is drawn.
